\chapter{Анализ технического, программного и информационного обеспечений}\label{ch:3}

\section{Анализ технического обеспечения}\label{sec:3.1}

Техническое обеспечение процесса предварительной оценки повреждений автомобилей по сути отсутствует. 

\begin{itemize}
    \item Клиентские устройства. Фотографии повреждений поступают преимущественно с мобильных устройств клиентов (смартфонов) через мессенджеры. Сотрудники автосервиса также используют мобильные устройства и/или персональные компьютеры для просмотра изображений и подготовки ответов.
    \item Сетевое подключение. Обмен изображениями осуществляется через Интернет и инфраструктуру мессенджеров. Специально выделенных сетей или VPN нет.
\end{itemize}

На текущем этапе нет каких-либо хранилищ или баз данных для систематического хранения фото и прочих метаданных: все изображения и информация о клиенте остаются в чатах мессенджеров и/или локально на устройствах сотрудников.
Резервное копирование не проводится.

Отсутствие какой-либо стандартизированной инфраструктуры приводит к невозможности систематизировать данные, к трудностям в аналитике и масштабировании, низкой отказоустойчивости процесса.

\section{Анализ программного обеспечения}\label{sec:3.2}

В текущем процессе отсутствует специализированное программное обеспечение. Клиенты используют только набор приложений без специализированного ПО для учёта, обработки и систематизации заявок.
Мессенджеры - основной канал приёма и обмена фотографиями. Клиенты присылают фото в личные чаты, сотрудники автосервиса дают ответы и проводят согласование.

\section{Анализ информационного обеспечения}\label{sec:3.3}

Основной источник информации — фотографии, присылаемые клиентами в мессенджеры. Текстовые сообщения с описанием повреждений отправляются клиентами в те же чаты.
Отсутствует какая-либо стандартизация формата присылаемых данных - изображения/описание повреждения могут прислать в любом виде, что и зачастую приводит к проблеме недостатка информации и трате времени на дополнительное обсуждение деталей.
Информация о состоянии обращения также нигде не хранится и сообщается клиенту лишь при наступлении определённого события, например, при завершении ремонта. 
